\section*{الفصل \ref{ch:b2:lab-techniques-2} --- تقنيات المختبر 2: التخليق والتحليل عمليًا}

\begin{solution}{exo:b2:lab-techniques-2:1}
جليد وماء؛ جليد وملح؛ جليد جاف وأسيتون (كما في أيونات الإينولات في \cref{ch:b2:enolates-aldol}).
\end{solution}

\begin{solution}{exo:b2:lab-techniques-2:2}
الطبقة السفلى. أضف بضع قطرات من الماء: تنضم إلى الطبقة المائية، وهي هنا العليا.
\end{solution}

\begin{solution}{exo:b2:lab-techniques-2:3}
تتكتل الدفعات الأولى وهي تمتص الماء؛ وحين يبقى الصلب المضاف حديثًا مفككًا يدور بحرية،
يكون المحلول جافًا.
\end{solution}

\begin{solution}{exo:b2:lab-techniques-2:4}
قيمة $R_f$ مرتفعة جدًا من أجل عمود (المطلوب نحو 0.3) والبقعتان متقاربتان: استعمل مُزيحًا أقل
قطبية، فيه مزيد من الهكسان (9\,:\,1 مثلًا)، مما يخفض قيمتي $R_f$ ويوسع الفصل بينهما مقيسًا
بحجوم العمود.
\end{solution}

\begin{solution}{exo:b2:lab-techniques-2:5}
\begin{itemize}
  \item هيدروجين كربونات الصوديوم (pH 8.3 $>$ 4.2 + 2): البنزوات في الماء؛ حمّض، واستخلص حمض البنزويك.
  \item هيدروكسيد الصوديوم (pH $>$ 12 $>$ 9.99 + 2): الفينوكسيد في الماء؛ حمّض، واستخلص الفينول.
  \item حمض الهيدروكلوريك (pH نحو 1 $<$ 4.6 $-$ 2): الأنيلينيوم في الماء؛ اجعل الوسط قاعديًا، واستخلص الأنيلين.
  \item يبقى النفثالين في الإيثر: جفف، وبخّر.
\end{itemize}
يجب أن يأتي الغسل بهيدروجين الكربونات أولًا: فالهيدروكسيد كان سيستخلص الحمض والفينول معًا.
% ledger: ph:NaHCO3, pka:benzoic, pka:phenol, pka:anilinium
\end{solution}

\begin{solution}{exo:b2:lab-techniques-2:6}
$1/T = 1/T_b - (R/\Delta_{\mathrm{vap}}H)\ln(p/p^\circ)$: البروموبنزين نحو \qty{49}{\degreeCelsius}،
والبنزالدهيد نحو \qty{68}{\degreeCelsius} (شكل \cref{prop:b2:lab-techniques-2:reduced-pressure}).
\end{solution}

\begin{solution}{exo:b2:lab-techniques-2:7}
$0.050 \times 200\,000/400 = \qty{25}{K}$.
\end{solution}

\begin{solution}{exo:b2:lab-techniques-2:8}
لكل مركب استجابته الخاصة: فالنسبة المئوية للمساحة لا تكون نسبة كتلية إلا إذا تساوت معاملات
الاستجابة. والرنين المغناطيسي النووي: $0.12/3 = 0.040$~mol من الشائبة لكل مول من الناتج، أي $0.040 \times 100 = \qty{4}{g}$
لكل \qty{200}{g} من الناتج؛ $w = 200/204 \approx 98.0~\%$.
\end{solution}

\begin{solution}{exo:b2:lab-techniques-2:9}
الكتلة النظرية $0.0300 \times 164 = \qty{4.92}{g}$. غير المصحح $4.10/4.92 \approx 83.3~\%$؛ والمصحح
$0.950 \times 4.10/4.92 \approx 79.2~\%$.
\end{solution}

\begin{solution}{exo:b2:lab-techniques-2:10}
يعني البنزين أن الكاشف تكوّن ثم بُرتن: ماء في الأدوات الزجاجية أو المذيب أو
الألدهيد؛ أو هواء (رطوبة) دخل عبر تسرب؛ أو ألدهيد تأكسد جزئيًا إلى حمض البنزويك، الذي يتلف
هيدروجينه الحمضي الكاشف. والعلاج: أدوات زجاجية مجففة في الفرن، ومذيب جاف، وألدهيد مقطَّر حديثًا،
وضغط زائد من النيتروجين يُتحقق منه عند الفقاعة. أما قلة الكاشف أصلًا فتدل على
إخفاق البدء (سطح مغنيزيوم مغطى بالأكسيد)، ويُتجنب ذلك ببُرادة حديثة أو منشَّطة.
\end{solution}

\begin{solution}{exo:b2:lab-techniques-2:11}
\qty{270}{kJ} عند \qty{150}{W}: على الأقل $270\,000/150 = \qty{1800}{s}$، أي نصف ساعة، وبشرط أن
يتفاعل الهاليد بالسرعة التي يُضاف بها. وإذا كانت الإضافة أسرع، تجاوز إطلاق الحرارة التبريد: فترتفع درجة الحرارة،
أو يتراكم الكاشف وقد يتفاعل لاحقًا دفعة واحدة (\cref{prop:b2:lab-techniques-2:accumulation}).
\end{solution}

\begin{solution}{exo:b2:lab-techniques-2:12}
إعادة التبلور: $13.0 \times 0.80 = \qty{10.4}{g}$ بنقاوة 99~\%. والعمود: $13.0 \times 0.95 =
\qty{12.4}{g}$ بنقاوة 99.5~\%، لكن مع \qty{1.5}{L} من المذيب يُشترى ويُبخَّر ويُتخلص منه. يعطي العمود
نحو \qty{2}{g} زيادة من ناتج أنقى قليلًا؛ وإعادة التبلور أبسط وأرخص وأقل
هدرًا، وتُفضَّل حين تكفي نقاوتها.
\end{solution}

\begin{solution}{pb:b2:lab-techniques-2:1}
\textbf{1.} الإيثوكسي إيثان: شديد الاشتعال، ضار، يكوّن فوق أكاسيد. والبروموبنزين: قابل للاشتعال، سام،
سام للأحياء المائية. والمغنيزيوم: صلب قابل للاشتعال، يطلق الهيدروجين مع الماء.
\textbf{2.} \ce{C6H5MgBr + H2O -> C6H6 + Mg(OH)Br}: كل قطرة ماء تتلف شيئًا من الكاشف؛ والتجفيف في الفرن
يزيل غشاء الماء عن الزجاج.
\textbf{3.} لإبعاد الرطوبة والأكسجين، الذي يتلف الكاشف أيضًا.
\textbf{4.} تدفق النيتروجين؛ وتمنع الهواء من الدخول عبر المخرج.
\textbf{5.} لا يتفاعل مع الكاشف، وأزواج أكسجينه الحرة تثبّت المغنيزيوم، 
ودرجة غليانه المنخفضة (\qty{307.7}{K}) تتيح للارتداد أن يحد درجة الحرارة.
\textbf{6.} $15.7/157.01 = \qty{0.1000}{mol}$ من البروموبنزين؛ $2.67/24.305 = \qty{0.1099}{mol}$ من المغنيزيوم،
بفائض 10~\%.
\textbf{7.} $0.1000 \times 270 = \qty{27.0}{kJ}$.
\textbf{8.} كي تنطلق الحرارة بالمعدل الذي يزيلها به الارتداد والحمام.
\textbf{9.} $0.0200 \times 270 = \qty{5.40}{kJ}$.
\textbf{10.} $5400/170 \approx \qty{32}{K}$.
\textbf{11.} نحو \qty{52}{\degreeCelsius}، أعلى بكثير من درجة غليان الإيثوكسي إيثان
(\qty{34.6}{\degreeCelsius}): غليان عنيف قد يغمر المكثف ويقذف بخارًا قابلًا للاشتعال
خارج الدورق.
\textbf{12.} أضف نحو عُشر البروموبنزين، وانتظر بدء التفاعل (تعكّر، وارتداد
لطيف)، ثم أضف الباقي بمعدل يحفظ ارتدادًا لطيفًا، مع إبقاء حمام الجليد جاهزًا.
\textbf{13.} \ce{C6H5MgBr + C6H5CHO} يعطي الألكوكسيد \ce{(C6H5)2CHOMgBr}؛ والبنزالدهيد،
$9.55/106.12 = \qty{0.0900}{mol}$، هو الكاشف المحدد.
\textbf{14.} الكحول بنزيلي مرتين: فالحمض القوي كان سيحوّله إلى كاتيون مستقر،
يقود إلى إيثرات أو نواتج استبدال؛ أما كلوريد الأمونيوم فحمضي بقدر يكفي لبرتنة الألكوكسيد
وإذابة أملاح المغنيزيوم.
\textbf{15.} الطبقة العليا: فالإيثوكسي إيثان أقل كثافة من الماء.
\textbf{16.} يزيل معظم الماء المذاب في الإيثر.
\textbf{17.} كبريتات المغنيزيوم حتى يصير حرّ الانسياب، فالترشيح، ثم مبخر دوّار تحت ضغط
مخفَّض.
\textbf{18.} يقترن بروميد فينيل المغنيزيوم بالبروموبنزين غير المتفاعل؛ ويحابيه تركيز موضعي مرتفع من
البروموبنزين ودرجة حرارة مرتفعة، وهذا سبب آخر لإضافة بطيئة.
\textbf{19.} ثنائي الفينيل ($R_f$ 0.85)، الأقل قطبية.
\textbf{20.} نحو $1/0.85 \approx 1.2$ حجم عمودي لثنائي الفينيل، و$1/0.25 = 4$ للكحول.
\textbf{21.} يسهم الكحول بعشرة بروتونات \omterm{def:b2:huckel:aromatic}{عطرية} لكل CH؛ والزيادة $10.90 - 10.00 = 0.90$ تأتي من
ثنائي الفينيل، بعشرة بروتونات لكل جزيء: 0.090 mol من ثنائي الفينيل لكل مول من الكحول.
\textbf{22.} $w = 184.24/(184.24 + 0.090 \times 154.21) \approx 0.930$.
\textbf{23.} $12.50 \times 0.930 \approx \qty{11.62}{g}$.
\textbf{24.} $0.0900 \times 184.24 \approx \qty{16.58}{g}$.
\textbf{25.} $12.50/16.58 \approx 75.4~\%$.
\textbf{26.} $11.62/16.58 \approx \boldsymbol{70~\%}$ (70.1~\%)، مصححًا بالنقاوة.
\end{solution}
